\section{Atingimento dos objetivos do trabalho}

O objetivo geral deste trabalho foi desenvolver e analisar um modelo dinâmico integrado entre a dinâmica orbital relativa e a dinâmica acoplada de uma \gls{etmr}, aplicado às fases de aproximação de longa distância, curta distância e aproximação final das missões de \gls{rvdb}, com validação em ambiente de simulação.

O objetivo do trabalho foi atingido.
O modelo desenvolvido permitiu representar o movimento relativo nos referenciais \gls{eci} e \gls{lvlh}, bem como a dinâmica da espaçonave perseguidora tanto na condição de corpo rígido, durante a fase não-operacional, como na condição acoplada ao \gls{mr}, por meio de formulação lagrangiana baseada na energia cinética do conjunto.

As simulações demonstraram a execução das fases de aproximação de longa distância, aproximação de curta distância e fase final, com correção e mantenimento de atitude através da lei de controle \gls{pd}.
A movimentação do \gls{mr} gerou torques de reação com a formulação dinâmica, e a ferramenta foi inserida e mantida na interface de acoplamento dentro dos limites estabelecidos nas simulações.